% =====================================================================
\chapter{Funciones reales elementales y sus gráficas}
\label{tema:1}
% =====================================================================

Todo el cálculo trata de funciones reales de una variable real. Este tema fija el vocabulario (dominio, imagen, gráfica, monotonía, paridad, inversa), repasa el catálogo de funciones elementales que se usarán como ejemplos en el resto del módulo y presenta las funciones de activación de las redes neuronales como funciones elementales más. Se siguen las §1.1--1.5 de Stewart~\cite{stewart}, las §1.1--1.5 de OpenStax~\cite{osc1} y los capítulos~3 y~4 de Spivak~\cite{spivak}; la parte conjuntista (inyectividad, composición e inversa) se estudió en el módulo~0.2.

\section{Funciones reales de variable real}

\begin{definicion}[Función real, dominio, imagen y gráfica]\label{def:funcion}
Una \emph{función real de variable real} es una función $f\colon D\to\R$ con $D\subseteq\R$. El conjunto $D=\dom f$ es su \emph{dominio}; $f(D)=\{f(x):x\in D\}$ es su \emph{imagen} o recorrido, y su \emph{gráfica} es el conjunto $\{(x,f(x)):x\in D\}\subseteq\R^{2}$. Cuando una función se da por una fórmula sin indicar el dominio, se sobrentiende el \emph{dominio natural}: el mayor conjunto de reales en el que la fórmula tiene sentido.
\end{definicion}

\begin{ejemplo}[Dominios naturales]\label{ej:dominios}
(a) $f(x)=\sqrt{4-x^{2}}$: se necesita $4-x^{2}\ge0$, luego $\dom f=[-2,2]$ y la imagen es $[0,2]$ (semicircunferencia superior de radio~2). (b) $g(x)=\ln(x-1)/(x-3)$: se necesita $x>1$ y $x\ne3$, luego $\dom g=(1,3)\cup(3,\infty)$. (c) $h(x)=1/(1+e^{-x})$: el denominador es siempre mayor que~1, luego $\dom h=\R$.
\end{ejemplo}

Las funciones se suman, se multiplican y se dividen punto a punto, en la intersección de los dominios (y, para el cociente, donde el denominador no se anula). La \emph{composición} $(g\circ f)(x)=g(f(x))$ está definida en $\{x\in\dom f: f(x)\in\dom g\}$.

\begin{definicion}[Acotación, monotonía, paridad y periodicidad]\label{def:propiedades}
Sea $f\colon D\to\R$.
\begin{enumerate}
  \item $f$ está \emph{acotada superiormente} si existe $M$ con $f(x)\le M$ para todo $x\in D$; \emph{inferiormente}, si existe $m$ con $f(x)\ge m$; \emph{acotada}, si lo está de ambas formas, es decir, si existe $K$ con $\abs{f(x)}\le K$ para todo $x\in D$.
  \item $f$ es \emph{creciente} si $x<y\Rightarrow f(x)\le f(y)$, y \emph{estrictamente creciente} si $x<y\Rightarrow f(x)<f(y)$, para todos $x,y\in D$. Análogamente, (estrictamente) \emph{decreciente}. Se llama \emph{monótona} a una función creciente o decreciente.
  \item Si $D$ es simétrico ($x\in D\Rightarrow -x\in D$), $f$ es \emph{par} si $f(-x)=f(x)$ e \emph{impar} si $f(-x)=-f(x)$ para todo $x\in D$.
  \item $f$ es \emph{periódica} de periodo $T>0$ si $x+T\in D$ y $f(x+T)=f(x)$ para todo $x\in D$.
\end{enumerate}
\end{definicion}

La gráfica de una función par es simétrica respecto del eje de ordenadas, y la de una impar, respecto del origen. Una función impar definida en~$0$ cumple $f(0)=-f(0)$, luego $f(0)=0$.

\begin{proposicion}[Transformaciones de gráficas]\label{prop:transformaciones}
Sea $f$ una función, $c\in\R$ y $k>0$. La gráfica de
\begin{enumerate}
  \item $y=f(x)+c$ es la de $f$ trasladada $c$ unidades en vertical;
  \item $y=f(x-c)$ es la de $f$ trasladada $c$ unidades en horizontal (hacia la derecha si $c>0$);
  \item $y=kf(x)$ es la de $f$ dilatada verticalmente por el factor $k$;
  \item $y=f(kx)$ es la de $f$ comprimida horizontalmente por el factor $k$ (dilatada si $k<1$);
  \item $y=-f(x)$ e $y=f(-x)$ son sus simétricas respecto del eje de abscisas y del de ordenadas.
\end{enumerate}
\end{proposicion}

\begin{proof}
Se prueba (ii); las demás son análogas. $(x,y)$ está en la gráfica de $y=f(x-c)$ si y solo si $y=f(x-c)$, es decir, si y solo si $(x-c,y)$ está en la gráfica de $f$. Luego la nueva gráfica se obtiene sumando $c$ a la abscisa de cada punto de la gráfica de~$f$.
\end{proof}

\section{El catálogo de funciones elementales}

\begin{definicion}[Funciones elementales básicas]\label{def:elementales}
Se llaman \emph{elementales} las funciones que se obtienen, mediante sumas, productos, cocientes, composiciones e inversas, a partir de:
\begin{enumerate}
  \item los \emph{polinomios} $p(x)=a_nx^{n}+\dots+a_1x+a_0$ y las \emph{funciones racionales} $p/q$ (cocientes de polinomios, en los puntos donde $q\ne0$);
  \item las \emph{potencias} $x^{a}$ ($x>0$, $a\in\R$) y las raíces $\sqrt[n]{x}$;
  \item la \emph{exponencial} $\exp(x)=e^{x}$, biyección creciente de $\R$ en $(0,\infty)$, y su inversa, el \emph{logaritmo} $\ln\colon(0,\infty)\to\R$;
  \item las funciones trigonométricas $\sin$, $\cos$, $\tan$ y sus inversas $\arcsin$, $\arccos$, $\arctan$ (módulo~0.4).
\end{enumerate}
\end{definicion}

Se recuerdan las propiedades de la exponencial y del logaritmo (módulo~0.3): $e^{x+y}=e^{x}e^{y}$, $e^{0}=1$, $\ln(xy)=\ln x+\ln y$, $\ln(x^{a})=a\ln x$, $e^{\ln x}=x$ para $x>0$ y $\ln(e^{x})=x$ para todo $x$; además $a^{x}=e^{x\ln a}$ para $a>0$. El valor absoluto $\abs{x}=\max(x,-x)$ y la parte entera $\floor{x}$ (el mayor entero $\le x$) completan el repertorio de ejemplos.

\begin{definicion}[Funciones hiperbólicas]\label{def:hiperbolicas}
Para $x\in\R$ se definen
\[
\sinh x=\frac{e^{x}-e^{-x}}{2},\qquad \cosh x=\frac{e^{x}+e^{-x}}{2},\qquad \tanh x=\frac{\sinh x}{\cosh x}=\frac{e^{x}-e^{-x}}{e^{x}+e^{-x}}.
\]
\end{definicion}

\begin{proposicion}[Propiedades de las funciones hiperbólicas]\label{prop:hiperbolicas}
Para todo $x\in\R$:
\begin{enumerate}
  \item $\cosh$ es par, $\cosh x\ge1$, y $\sinh$ y $\tanh$ son impares.
  \item $\cosh^{2}x-\sinh^{2}x=1$ para todo $x$.
  \item $\cosh x+\sinh x=e^{x}$ y $\cosh x-\sinh x=e^{-x}$.
  \item $\abs{\tanh x}<1$, y $\tanh$ es estrictamente creciente de $\R$ en $(-1,1)$.
\end{enumerate}
\end{proposicion}

\begin{proof}
(i) Al cambiar $x$ por $-x$ se intercambian $e^{x}$ y $e^{-x}$; la suma no cambia y la diferencia cambia de signo. Por la desigualdad entre las medias aritmética y geométrica, $\cosh x\ge\sqrt{e^{x}e^{-x}}=1$. (iii) Es inmediato sumando y restando las definiciones. (ii) Por (iii), $\cosh^{2}x-\sinh^{2}x=(\cosh x+\sinh x)(\cosh x-\sinh x)=e^{x}e^{-x}=1$. (iv) Multiplicando numerador y denominador por $e^{x}$, $\tanh x=\dfrac{e^{2x}-1}{e^{2x}+1}=1-\dfrac{2}{e^{2x}+1}$. Como $e^{2x}+1>1$, se tiene $-1<\tanh x<1$; y como $x\mapsto e^{2x}+1$ es estrictamente creciente y positiva, $2/(e^{2x}+1)$ es estrictamente decreciente y $\tanh$ estrictamente creciente. Que la imagen es todo $(-1,1)$ se ve resolviendo $\tanh x=y$ (\cref{eje:1-arsinh}).
\end{proof}

El nombre procede de (ii): el punto $(\cosh t,\sinh t)$ recorre la rama derecha de la hipérbola $x^{2}-y^{2}=1$, como $(\cos t,\sin t)$ recorre la circunferencia $x^{2}+y^{2}=1$.

\section{Funciones de activación}

Una neurona artificial calcula una combinación lineal de sus entradas y le aplica una función no lineal fija, la \emph{función de activación}. Las más usadas son funciones elementales (o casi) y sirven como ejemplos recurrentes en este módulo.

\begin{definicion}[Funciones de activación]\label{def:activaciones}
Para $x\in\R$ se definen
\begin{enumerate}
  \item la \emph{sigmoide} o función logística: $\sigma(x)=\dfrac{1}{1+e^{-x}}$;
  \item la \emph{tangente hiperbólica} $\tanh x$ (\cref{def:hiperbolicas});
  \item la \emph{unidad lineal rectificada}: $\ReLU(x)=\max(0,x)$;
  \item la \emph{softplus}: $\softplus(x)=\ln(1+e^{x})$.
\end{enumerate}
\end{definicion}

\begin{figure}[h]
\centering
\begin{tikzpicture}[x=0.75cm,y=1.1cm]
  \foreach \y in {-1,1} {\draw[dashed,color=regla] (-4.2,\y)--(-0.55,\y); \draw[dashed,color=regla] (0.1,\y)--(4.2,\y);}
  \draw[->,color=regla] (-4.3,0)--(4.4,0) node[right,color=tinta]{\small $x$};
  \draw[->,color=regla] (0,-1.25)--(0,2.6);
  \foreach \t in {-4,-2,2,4} \draw[color=regla] (\t,0.04)--(\t,-0.04) node[below,color=tinta]{\scriptsize $\t$};
  \foreach \t in {-1,1,2} \draw[color=regla] (0.06,\t)--(-0.06,\t) node[left,color=tinta]{\scriptsize $\t$};
  \draw[thick,domain=-4.2:4.2,samples=80] plot (\x,{1/(1+exp(-\x))});
  \draw[thick,densely dashed,domain=-4.2:4.2,samples=80] plot (\x,{(exp(\x)-exp(-\x))/(exp(\x)+exp(-\x))});
  \draw[very thick,color=enlace] (-4.2,0)--(0,0)--(2.5,2.5);
  \draw[thick,densely dotted,domain=-4.2:2.0,samples=60] plot (\x,{ln(1+exp(\x))});
  \node[right] at (4.2,1.0) {\small $\sigma$};
  \node[left] at (-4.2,-1.0) {\small $\tanh$};
  \node[right,color=enlace] at (2.5,2.5) {\small $\ReLU$};
  \node[left] at (1.9,2.15) {\small $\softplus$};
\end{tikzpicture}
\caption{Gráficas de $\sigma$ (trazo continuo), $\tanh$ (discontinuo), $\ReLU$ (trazo grueso) y $\softplus$ (punteado).}\label{fig:activaciones}
\end{figure}

\begin{proposicion}[Propiedades de la sigmoide]\label{prop:sigmoide}
La función sigmoide cumple:
\begin{enumerate}
  \item $0<\sigma(x)<1$ para todo $x$, y $\sigma(0)=\tfrac12$.
  \item $\sigma(-x)=1-\sigma(x)$; en consecuencia, la gráfica es simétrica respecto del punto $(0,\tfrac12)$.
  \item $\sigma$ es estrictamente creciente y biyectiva de $\R$ en $(0,1)$, con inversa $\logit(p)=\ln\dfrac{p}{1-p}$.
  \item $\tanh x=2\sigma(2x)-1$.
\end{enumerate}
\end{proposicion}

\begin{proof}
(i) $1+e^{-x}>1$. (iii) $x\mapsto 1+e^{-x}$ es estrictamente decreciente y positiva, luego su inversa multiplicativa es estrictamente creciente. Dado $p\in(0,1)$, la ecuación $\sigma(x)=p$ equivale a $e^{-x}=\frac1p-1=\frac{1-p}{p}>0$, que tiene la solución única $x=\ln\frac{p}{1-p}$. Las demostraciones de (ii) y (iv) son el \cref{eje:1-sigmoide-simetria}.
\end{proof}

La propiedad (iv) dice que $\tanh$ es una sigmoide reescalada: misma forma, imagen $(-1,1)$ en lugar de $(0,1)$ y pendiente cuatro veces mayor en el origen (\cref{tema:4}). En la regresión logística (módulo~8), $\sigma$ transforma una puntuación real en una probabilidad, y $\logit$ deshace esa transformación: el logit de una probabilidad es el logaritmo de la razón de probabilidades $p/(1-p)$, de donde procede el término \emph{logits} con que se designan las salidas de una red antes de la softmax.

\begin{proposicion}[ReLU y softplus]\label{prop:relu}
Para todo $x\in\R$:
\begin{enumerate}
  \item $\ReLU(x)=\dfrac{x+\abs{x}}{2}$ y $\ReLU(x)-\ReLU(-x)=x$;
  \item $\softplus(x)-\softplus(-x)=x$;
  \item $\ReLU(x)<\softplus(x)\le\ReLU(x)+\ln 2$, con igualdad a la derecha solo en $x=0$.
\end{enumerate}
\end{proposicion}

\begin{proof}
(i) Si $x\ge0$, $(x+\abs{x})/2=x$; si $x<0$, vale~$0$. La segunda igualdad se sigue de la primera, porque $\abs{-x}=\abs{x}$. (ii) $\softplus(x)-\softplus(-x)=\ln\dfrac{1+e^{x}}{1+e^{-x}}=\ln\dfrac{e^{x}(e^{-x}+1)}{1+e^{-x}}=\ln e^{x}=x$. (iii) Si $x\le0$: $\softplus(x)=\ln(1+e^{x})$ con $0<e^{x}\le1$, luego $0<\softplus(x)\le\ln2$, con igualdad solo si $x=0$. Si $x>0$: por (ii), $\softplus(x)=x+\softplus(-x)$ y $0<\softplus(-x)<\ln 2$.
\end{proof}

La softplus es, por tanto, una versión «suavizada» de la ReLU que nunca se separa de ella más de $\ln 2\approx0{,}693$ (\cref{fig:activaciones}). Otras activaciones modernas, como la GELU, $\GELU(x)=x\,\Phi(x)$ con $\Phi$ la función de distribución de la normal estándar, no son elementales; se estudian en el módulo~9.

\section{Función inversa}

\begin{proposicion}[Monotonía estricta e inversa]\label{prop:inversa-monotona}
Sea $f\colon D\to\R$ estrictamente creciente. Entonces $f$ es inyectiva, su inversa $f^{-1}\colon f(D)\to D$ es estrictamente creciente, y la gráfica de $f^{-1}$ es la simétrica de la de $f$ respecto de la recta $y=x$.
\end{proposicion}

\begin{proof}
Si $x\ne y$, por ejemplo $x<y$, entonces $f(x)<f(y)$ y, en particular, $f(x)\ne f(y)$: $f$ es inyectiva y, por tanto, biyectiva sobre su imagen. Sean $u<v$ en $f(D)$, con $u=f(x)$ y $v=f(y)$. Si fuera $x\ge y$, el crecimiento de $f$ daría $u=f(x)\ge f(y)=v$, contradicción; luego $f^{-1}(u)=x<y=f^{-1}(v)$. Por último, $(a,b)$ está en la gráfica de $f$ si y solo si $b=f(a)$, si y solo si $a=f^{-1}(b)$, si y solo si $(b,a)$ está en la gráfica de $f^{-1}$; y $(a,b)\mapsto(b,a)$ es la simetría respecto de $y=x$.
\end{proof}

\begin{ejemplo}[Inversas de funciones elementales]\label{ej:inversas}
(a) $\exp\colon\R\to(0,\infty)$ y $\ln$ son inversas, ambas estrictamente crecientes. (b) $x\mapsto x^{2}$ no es inyectiva en $\R$, pero su restricción a $[0,\infty)$ es estrictamente creciente y su inversa es $\sqrt{\phantom{x}}$. (c) $\sigma$ y $\logit$ son inversas (\cref{prop:sigmoide}). (d) Para invertir $y=\frac{2x+1}{x-3}$ ($x\ne3$) se despeja: $y(x-3)=2x+1$, $x(y-2)=3y+1$, $x=\frac{3y+1}{y-2}$, válido para $y\ne2$; la imagen es $\R\setminus\{2\}$.
\end{ejemplo}

\begin{observacion}[Errores frecuentes]\label{obs:1-errores}
(a) Confundir $f^{-1}(x)$ (función inversa) con $1/f(x)=f(x)^{-1}$ (inverso multiplicativo): $\sin^{-1}$ suele denotar $\arcsin$, no $1/\sin$.
(b) Usar $\ln(a+b)=\ln a+\ln b$ o $e^{a+b}=e^{a}+e^{b}$, que son falsas; las correctas son $\ln(ab)=\ln a+\ln b$ y $e^{a+b}=e^{a}e^{b}$.
(c) Olvidar el dominio al componer: $\ln(\ln x)$ exige $x>1$, no solo $x>0$.
(d) Pensar que una función que no es par es impar: la mayoría no son ni lo uno ni lo otro (por ejemplo, $e^{x}$), aunque toda función de dominio simétrico es suma de una par y una impar (\cref{eje:1-par-impar}).
(e) En $y=f(x-c)$, trasladar la gráfica hacia la izquierda cuando $c>0$; el desplazamiento es hacia la derecha.
\end{observacion}

\begin{observacion}[Conexión con el curso]
La sigmoide y la tangente hiperbólica fueron las activaciones clásicas de las redes neuronales; su saturación (para $\abs{x}$ grande se aplanan, \cref{fig:activaciones}) provoca los gradientes que se desvanecen del módulo~11. La ReLU, que no satura para $x>0$, y sus variantes suaves (softplus, GELU, SiLU $x\sigma(x)$) son las activaciones de los perceptrones multicapa del módulo~9 y de los bloques del Transformer de los módulos~14 y~18. La función $\logit=\sigma^{-1}$ explica el nombre de las salidas de un LLM antes de la softmax (módulo~1).
\end{observacion}

\begin{lecturas}
\item Stewart~\cite{stewart}, §1.1 «Four Ways to Represent a Function», §1.3 «New Functions from Old Functions» y §1.5 «Inverse Functions and Logarithms».
\item OpenStax~\cite{osc1}, §1.1 «Review of Functions», §1.2 «Basic Classes of Functions», §1.4 «Inverse Functions» y §1.5 «Exponential and Logarithmic Functions» (incluye las funciones hiperbólicas).
\item Spivak~\cite{spivak}, cap.~3 «Functions» y cap.~4 «Graphs».
\item MIT 18.01SC~\cite{mit1801}, sesión~20 «Hyperbolic Trig Functions».
\end{lecturas}

\section*{Ejercicios del tema 1}
\addcontentsline{toc}{section}{Ejercicios del tema 1}

\begin{ejercicio}{1}{eje:1-dominios}
Determine el dominio natural de: (a) $f(x)=\sqrt{x^{2}-5x+6}$; (b) $g(x)=\ln\dfrac{1+x}{1-x}$; (c) $h(x)=\dfrac{1}{\ln x}$; (d) $k(x)=\sqrt{\ln x}$.
\end{ejercicio}

\begin{ejercicio}{1}{eje:1-transformaciones}
La gráfica de $f$ pasa por los puntos $(0,0)$, $(1,2)$ y $(3,-1)$. ¿Por qué puntos pasa la gráfica de: (a) $y=f(x-2)+1$; (b) $y=3f(x)$; (c) $y=f(2x)$; (d) $y=-f(-x)$? Describa en cada caso la transformación con la \cref{prop:transformaciones}.
\end{ejercicio}

\begin{ejercicio}{2}{eje:1-paridad}
Estudie si son pares, impares o ninguna de las dos cosas: (a) $x^{3}-x$; (b) $\abs{x}+\cos x$; (c) $e^{x}$; (d) $\ln\dfrac{1-x}{1+x}$ en $(-1,1)$; (e) $x\,\sigma(x)$; (f) $\softplus(x)-\dfrac{x}{2}$.
\end{ejercicio}

\begin{ejercicio}{2}{eje:1-composicion}
Sean $f(x)=\sqrt{x}$ y $g(x)=1-x^{2}$. Calcule $f\circ g$, $g\circ f$, $f\circ f$ y $g\circ g$, y el dominio de cada una.
\end{ejercicio}

\begin{ejercicio}{2}{eje:1-ecuaciones}
Resuelva las ecuaciones: (a) $e^{3x}-15=0$; (b) $7^{3x-2}=11$; (c) $\ln x+\ln(x-2)=\ln 4$; (d) $4^{x+1}-32=0$.
\fuente{adaptado de OpenStax~\cite{osc1}, §1.5, ejercicios de la sección (ecuaciones exponenciales y logarítmicas).}
\end{ejercicio}

\begin{ejercicio}{2}{eje:1-hiperbolicas-calculo}
Escriba con exponenciales y simplifique: (a) $2\cosh(\ln x)$ para $x>0$; (b) $\cosh 4x+\sinh 4x$; (c) $\cosh 2x-\sinh 2x$; (d) $\ln(\cosh x+\sinh x)+\ln(\cosh x-\sinh x)$.
\fuente{adaptado de OpenStax~\cite{osc1}, §1.5, ejercicio de la sección sobre funciones hiperbólicas.}
\end{ejercicio}

\begin{ejercicio}{2}{eje:1-sigmoide-valores}
(a) Calcule $\sigma(0)$, $\sigma(\ln 3)$ y $\sigma(-\ln 3)$. (b) Resuelva $\sigma(x)=0{,}9$ y $\sigma(x)=0{,}01$. (c) ¿Para qué valores de $x$ es $\sigma(x)>0{,}99$? (d) Calcule $\tanh(\ln 2)$.
\end{ejercicio}

\begin{ejercicio}{2}{eje:1-inversa-racional}
Sea $f(x)=\dfrac{x+1}{x-1}$ en $\R\setminus\{1\}$. (a) Calcule $f\circ f$ y deduzca que $f$ es su propia inversa. (b) Halle la imagen de $f$. (c) Explique geométricamente por qué la gráfica de $f$ es simétrica respecto de la recta $y=x$.
\end{ejercicio}

\begin{ejercicio}{3}{eje:1-sigmoide-simetria}
Demuestre las propiedades (ii) y (iv) de la \cref{prop:sigmoide}: $\sigma(-x)=1-\sigma(x)$ y $\tanh x=2\sigma(2x)-1$ para todo $x\in\R$. Deduzca que $\sigma(x)=\tfrac12\bigl(1+\tanh\tfrac{x}{2}\bigr)$.
\end{ejercicio}

\begin{ejercicio}{3}{eje:1-hiperbolicas-suma}
Demuestre que, para todos $x,y\in\R$:
(a) $\sinh(x+y)=\sinh x\cosh y+\cosh x\sinh y$;
(b) $\cosh(x+y)=\cosh x\cosh y+\sinh x\sinh y$;
(c) $\tanh(x+y)=\dfrac{\tanh x+\tanh y}{1+\tanh x\tanh y}$.
Compare con las fórmulas trigonométricas del módulo~0.4.
\fuente{adaptado de MIT 18.01SC~\cite{mit1801}, sesión~20 «Hyperbolic Trig Functions».}
\end{ejercicio}

\begin{ejercicio}{3}{eje:1-arsinh}
(a) Demuestre que $\sinh\colon\R\to\R$ es biyectiva y que su inversa es $\arsinh y=\ln\bigl(y+\sqrt{y^{2}+1}\bigr)$. (b) Demuestre que $\tanh\colon\R\to(-1,1)$ es biyectiva y que $\artanh y=\tfrac12\ln\dfrac{1+y}{1-y}$. (c) Compruebe que $\artanh y=\tfrac12\logit\bigl(\tfrac{1+y}{2}\bigr)$.
\end{ejercicio}

\begin{ejercicio}{3}{eje:1-estricta}
Sea $f\colon[a,b]\to\R$ estrictamente creciente. Demuestre que $f$ es inyectiva y que, si $f(x)<f(y)$, entonces $x<y$. Muestre con un ejemplo que una función inyectiva en $[0,2]$ no tiene por qué ser monótona.
\fuente{adaptado de Cambridge, \emph{Tripos} IA 2013, Paper~1~\cite{tripos2013}, pregunta~10E(b)(i).}
\end{ejercicio}

\begin{ejercicio}{3}{eje:1-monotonia-composicion}
Sean $f$ y $g$ funciones de $\R$ en $\R$. Demuestre: (a) si $f$ y $g$ son crecientes, $g\circ f$ es creciente; (b) si $f$ es creciente y $g$ decreciente, $g\circ f$ es decreciente; (c) si $f$ y $g$ son crecientes, $f+g$ es creciente; (d) el producto de dos funciones crecientes no es necesariamente creciente. Deduzca de (a) y (b) que $x\mapsto\sigma(-x^{3})$ es decreciente.
\end{ejercicio}

\begin{ejercicio}{3}{eje:1-par-impar}
Sea $D\subseteq\R$ simétrico y $f\colon D\to\R$. Demuestre que existen una única función par $p$ y una única función impar $q$ con $f=p+q$, y obtenga fórmulas para ellas. Identifíquelas cuando $f(x)=e^{x}$ y cuando $f(x)=\softplus(x)$.
\end{ejercicio}

\begin{ejercicio}{3}{eje:1-softplus}
Demuestre que $\softplus(x)=x+\softplus(-x)$ y que $\softplus(x)=-\ln\sigma(-x)$ para todo $x$. Deduzca que $\softplus$ es estrictamente creciente y que $\ln\sigma(x)=-\softplus(-x)$. Explique por qué esta última identidad es útil para calcular $\ln\sigma(x)$ con $x$ muy negativo.
\end{ejercicio}
